\chapter{Number Theory}

\subsection{Modular Integer (Mint)}
\kactlimport{mint.h}

\subsection{nCr}
\kactlimport{nCr.h}

\subsection{Euclid}
\kactlimport{euclid.h}

\subsection{Sieve of Eratosthenes}
\kactlimport{Eratosthenes.h}

\subsection{some functions}
Let $n = p_1^{a_1} p_2^{a_2} \cdots p_k^{a_k}$.

\[
\begin{aligned}
\text{NOD: } \tau(n) &= \prod_{i=1}^k (a_i + 1) \\
\text{SOD: } \sigma(n) &= \prod_{i=1}^k \frac{p_i^{a_i+1} - 1}{p_i - 1} \\
\text{Euler: } \phi(n) &= n \prod_{i=1}^k \left(1 - \frac{1}{p_i}\right)
\end{aligned}
\]


\subsection{Euler Phi Function}
\kactlimport{phi.h}

\subsection{Sieve Factorization}
\kactlimport{Sievefactor.h}

\subsection{Linear Diophantine Equations}
\kactlimport{Lde.h}

\subsection{Chinese Remainder Theorem}
\kactlimport{crt.h}

\subsection{Mobius Function}
\kactlimport{mobius.h}

\subsection{Fast Factorization}
\kactlimport{fastfactor.h}

\section{Big Integer Utilities}

\subsection{String Addition}
\kactlimport{stringAdd.cpp}

\subsection{String Divisible Check}
\kactlimport{stringDivisible.cpp}

\subsection{String Multiplication}
\kactlimport{stringMul.cpp}

\subsection{safin's NT temp}


% %====================== Sieve and Prime Related ======================
% \section{Sieve and Prime Utilities}
% \kactlimport{ETFsieve.cpp}        % new added, last month
% \kactlimport{ETFsqrt.cpp}         % new added, last month
 \kactlimport{FactoriseSieve.cpp}  % new added, last month
% \kactlimport{FactoriseSqrt.cpp}   % added mod, last week
 \kactlimport{NODandSODsieve.cpp}  % added math pdf from kactl, last week
% \kactlimport{NodSqrt.cpp}         % new added, last month
% \kactlimport{PrimeSieve.cpp}      % new added, last month
% % \kactlimport{PrimeSqrt.cpp}       % new added, last month
 \kactlimport{SPFandHPF.cpp}       % spf add korlam, last week
 \kactlimport{SegmentedSieve.cpp}  % new added, last month
% \kactlimport{SodSqrt.cpp}         % new added, last month
 \kactlimport{bitsetSieve.cpp}     % memory efficient Sieve, last week
% \kactlimport{progkriyaSieve.cpp}  % eta boi er sieve ta, last week
 \kactlimport{sieve upto 1e9.cpp}  % source added, last week

% %====================== Modular Arithmetic ===========================
% \section{Modular Arithmetic}
\kactlimport{mod bin exp.cpp} 
% \kactlimport{modLog.cpp}
% \kactlimport{Mint.cpp}

% %====================== Combinatorics ================================
% \section{Combinatorics}
% \kactlimport{nCr nPr.cpp}


% %====================== Big Integer Utilities ========================
% \section{Big Integer Utilities}
% \kactlimport{stringAdd.cpp}
% \kactlimport{stringDivisible.cpp}
% \kactlimport{stringMul.cpp}

% %====================== Others ====================================
% \section{Others}
% \kactlimport{CRT.cpp}
% \kactlimport{fibonacci number faster.cpp}
% \kactlimport{EGCD.cpp}
% \kactlimport{EGCD LDE.cpp}
% 	% \subsection{Bézout's identity}
% 	% For $a \neq $, $b \neq 0$, then $d=gcd(a,b)$ is the smallest positive integer for which there are integer solutions to $ax+by=d$. If $(x,y)$ is one solution, then all solutions are given by $\left(x+kb/d, y-ka/d\right)$, $k\in\mathbb{Z}$. Find one solution using egcd.


\section{More Number Theory}

\[
\sum_{k=1}^n \frac{1}{\gcd(k,n)}
= \sum_{d\mid n} \frac{1}{d}\,\phi\!\left(\frac{n}{d}\right)
= \frac{1}{n} \sum_{d\mid n} d\,\phi(d)
\]

\[
\sum_{k=1}^n \frac{k}{\gcd(k,n)}
= \frac{n}{2} \cdot \frac{1}{n}
  \sum_{d\mid n} d\,\phi(d)
\]

\[
\sum_{k=1}^{n} \frac{n}{\gcd(k,n)}
= 2 \sum_{k=1}^{n} \frac{k}{\gcd(k,n)} - 1,
\quad n > 1
\]

\[
\sum_{i=1}^n \sum_{j=1}^n [\gcd(i,j)=1]
= \sum_{d=1}^n \mu(d)\left\lfloor \frac{n}{d} \right\rfloor^2
\]

\[
\sum_{i=1}^n \sum_{j=1}^n \gcd(i,j)
= \sum_{d=1}^n \phi(d)\left\lfloor \frac{n}{d} \right\rfloor^2
\]

\[
\sum_{i=1}^n \sum_{j=1}^n i\cdot j\,[\gcd(i,j)=1]
= \sum_{i=1}^n \phi(i)\,i^2
\]

\[
\sum_{i=1}^{n}\sum_{j=1}^{n}\mathrm{lcm}(i,j)
= \sum_{l=1}^n
\left(
\frac{(1+\lfloor n/l\rfloor)\lfloor n/l\rfloor}{2}
\right)^2
\]

\subsection{Permutations}
\subsubsection{Factorial Values}

\begin{center}
\begin{tabular}{c|cccccccc}
$n$  & 1 & 2 & 3 & 4 & 5 & 6 & 7 &8\\
\hline
$n!$ & 1 & 2 & 6 & 24 & 120 & 720 & 5040 &40320\\
\end{tabular}

\vspace{3mm}

\begin{tabular}{c|ccccc}
$n$  & 9 & 10 & 11 & 12 & 13\\
\hline
$n!$ & 362880 & 3628800 & 4.0e7 & 4.8e8 & 6.2e9\\
\end{tabular}

\vspace{3mm}

\begin{tabular}{c|ccccccc}
$n$  & 14 & 15 & 16 & 17 & 20 &25\\
\hline
$n!$ & 8.7e10 & 1.3e12 & 2.1e13 & 3.6e14 & 2e18 &2e25\\
\end{tabular}

\end{center}




% \section{Fractions}
% 	%\kactlimport{ContinuedFractions.h}
% 	\kactlimport{FracBinarySearch.h}

\section{Mobius Function}
\[
	\mu(n) = \begin{cases} 0 & n \textrm{ is not square free}\\ 1 & n \textrm{ has even number of prime factors}\\ -1 & n \textrm{ has odd number of prime factors}\\\end{cases}
\]
  Mobius Inversion:
  \[ g(n) = \sum_{d|n} f(d) \Leftrightarrow f(n) = \sum_{d|n} \mu(d)g(n/d) \]
  Other useful formulas/forms:

  $ \sum_{d | n} \mu(d) = [ n = 1] $,
  $ \phi(n) = \sum_{d | n} \mu(d)\frac{n}{d}$

  $ g(n) = \sum_{n|d} f(d) \Leftrightarrow f(n) = \sum_{n|d} \mu(\frac{d}{n})g(d)$

  $ g(n) = \sum_{1 \leq m \leq n} f(\left\lfloor\frac{n}{m}\right \rfloor ) \Leftrightarrow f(n) = \sum_{1\leq m\leq n} \mu(m)g(\left\lfloor\frac{n}{m}\right\rfloor)$

  If $f$ multiplicative, $\sum_{d | n} \mu(d)f(d) = \prod_{\textnormal{prime } p | n}(1 - f(p))$ and $\sum_{d | n}\mu^2(d)f(d) = \prod_{\textnormal{prime } p | n} (1 + f(p))$. 

  If $s_f(n) = \sum_{i = 1}^{n} f(i)$ is a prefix sum of mulitplicative $f$ then $s_{f * g}(n) = \sum_{1 \leq xy \leq n}f(x)g(y)$. Then $s_f(n) = \{s_{f * g}(n) - \sum_{d = 2}^{n} s_f(\lfloor n / d \rfloor) g(d)\} / g(1)$ where $f * g(n) = \sum_{d | n} f(d)g(n / d)$ (Dirichlet). Precompute (linear sieve) $O(n^{2/3})$ first values of $s_f$ for complexity $O(n^{2/3})$.

  Useful sums and convolutions: $\epsilon = \mu * \textnormal{\textbf{1}}$, id = $\phi * \textnormal{\textbf{1}}$, id = $g * \textnormal{id}_2$, where $\epsilon(n) = [n = 1]$, $\textnormal{\textbf{1}}(n) = 1$, id$(n) = n$, id$_k(n) = n^k$, $g(n) = \sum_{d | n}\mu(d)nd$.\\
  coprime pairs in $[1,n]$ is $\sum_{d = 1}^{n}\mu(d)\lfloor n / d \rfloor ^2$. Sum of GCD pairs in $[1, n]$ is $\sum_{d = 1}^{n}\phi(d)\lfloor n / d \rfloor ^2$. Sum of LCM pairs in $[1, n]$ is $\sum_{d = 1}^{n}(\frac{\lfloor n / d \rfloor (1 + \lfloor n / d \rfloor)}{2})^2 g(d)$, where $g$ is defined above with $g(p^k) = p^k - p^{k+1}$.

